% Generated by build/make_tables.py from calibration_error_by_year.csv. Do not edit.

\begin{table}[tbp]
\centering
\small
\caption{Sizing check: the disclosed sleeve targets, and the largest error the calibration leaves behind}
\label{tab:sizing}
\begin{tabular}{@{}lrrrrr@{}}
\toprule
 & \multicolumn{4}{c}{Disclosed target (\$mm)} & Largest error \\
\cmidrule(lr){2-5}
Year & Rates & Equity & Currency & Commodity & (\% of target) \\
\midrule
2016 & 45 & 25 & 21 & 17 & 4.26e-14 \\
2017 & 40 & 24 & 12 & 13 & 1.78e-14 \\
2018 & 46 & 31 & 14 & 11 & 1.61e-14 \\
2019 & 46 & 27 & 11 & 12 & 3.23e-14 \\
2020 & 71 & 55 & 23 & 20 & 3.09e-14 \\
2021 & 60 & 43 & 13 & 25 & 3.30e-14 \\
2022 & 97 & 33 & 32 & 47 & 3.02e-14 \\
2023 & 96 & 29 & 24 & 19 & 2.96e-14 \\
2024 & 81 & 37 & 26 & 19 & 3.84e-14 \\
2025 & 70 & 47 & 24 & 17 & 3.02e-14 \\
\bottomrule
\end{tabular}
\notetext{Targets are Goldman's disclosed average daily VaR by category, as first reported. The final column is the largest absolute sizing error across the four sleeves in that year, as a percentage of target. Across all 40 sleeve-years the largest error is 4.26e-14\%, which is floating-point noise: the calibration is exact, as Equation~\eqref{eq:homogeneity} requires, and this table verifies that claim rather than reporting a convergence result.}
\end{table}
